\section{Selmer conditions on the Taylor--Wiles limit}

Fix the single group system $T$ of R70P1. Assume $T.\mathrm{ProL}$,
so the relative pro-$\ell$ residual is trivial at the base and at
every finite level. Fix one of its normalized points $x$, its
package $k_x$, and an indexed list of distinct roots satisfying
$T.\mathrm{Eig}$. Set $A=K[[t]]$, $A_m=A/(t^m)$ and
$E=K((t))$. Write $g$ for the number of generators of $k_x$,
$d$ for the exponent parameter of $B1$, and $Q=|K|$. Assume
the native $B1(k_x,d)$ and $B2(k_x)$ for this same point. The
group of the limit maps to the original group by $h$.

% BEGIN L71C1
\paragraph{L71C1. The same finite Selmer cocycles at an ultrafilter-large set of levels.}
Let $P$ be the finitely generated profinite cover of LIM18, with
its surjections $\pi_N:P\to\Gamma_N$ and
$\pi_\infty:P\to\Gamma_\infty$ and its distinguished conjugation
lift. The two projections to the base agree modulo the subgroup
$T.H$ in the sign kernel, so these maps have the same sign.
Their restrictions on the
sign kernel $P^0$ are onto the respective index-two subgroups.
The group $P^0$ is topologically finitely generated by finite-index
Schreier generators. Precisely, choose a finite generating set
of a dense abstract subgroup of $P$. Its intersection with the
open subgroup $P^0$ is dense in $P^0$ and has index at most two
in the dense subgroup. Schreier's lemma applies there.

Fix $m>0$. Use any of the following actions over $A_m$: the
ordinary matrix action on even matrices, or the action of the
first-order representation on $W_1$ or $W_2$. In the latter
cases the coefficient algebra has basis
$1,e_j,f_j,\eta$ for $1\le j\le g$, with
$e_i f_j=f_j e_i=\delta_{ij}\eta$ and other nonconstant products
zero. The acting matrices $R=(1+\sum e_j\kappa_j)r$,
the flag $(1-\sum e_jY_{j,v})w_v$ and its inverse all reduce
from the fixed integral data. The modules are given by setting
specified entries in these finitely many coordinates to zero.
They are finite; the representations are continuous since
$r,\kappa_j$ are continuous at each precision. The limit's
convergence of the base projections implies that the actions and
matrix representations on $P^0$ at level $N$ agree with those
at infinity on an ultrafilter-large set. To see uniformity in
the group element, use the normal core in the full base group
of the kernel of the finite representation. The assertion
$\mathrm{p\_lim}$ holds simultaneously on the entire cover
modulo this open normal subgroup.

For the resulting fixed finite action there are finitely many
continuous crossed homomorphisms on $P^0$: the associated map
to the finite semidirect product of the module and the acting group
is a homomorphism determined by values at topological generators.
Choose one open normal subgroup $V\subset P^0$, normal in $P$,
killing the actions and cocycle values there for \emph{all}
these homomorphisms, by intersecting the kernels of their affine
maps and taking the normal core. Modulo $V$ the limit identifies
$\ker\pi_N$ with $\ker\pi_\infty$ (their products with $V$
coincide) at ultrafilter-many levels. Therefore a fixed cocycle
on the cover descends through one exactly when it descends
through the other. Continuity descends because the surjections
are continuous quotient maps of profinite groups.

The local and tame maps into the two groups are uniformly close
modulo $V$ on this same large set, after a finite intersection.
This means for each corresponding pair of local elements there
exist lifts to $P$ agreeing modulo $V$. For a local element
in the sign kernel both lifts lie in $P^0$. Hence the pulled-back cocycle
and representation take identical values on the two elements.
Conjugation uses the same lift in $P$. The block, transpose,
tame vanishing, upper-triangular and torsion-diagonal conditions
can thus be tested with the same local correcting matrices.
This gives equality of the two Selmer sets as functions on
$P^0$ under pullback, and equality of their boundary sets since
a skew primitive gives the same function under equal actions.
The finitely many choices of action just listed can all be treated
at once. At $m=0$ the relevant target modules are zero.
% END L71C1

% BEGIN L71C2
\paragraph{L71C2. Both finite bounds reach the limit with their original exponents.}
For each $m$ intersect the set of L71C1 with the eventual
$B1$ bounds. A cofinal ultrafilter contains each eventual set.
At any level in the intersection, pullback on the surjective
maps gives identical finite cardinalities. The bound at infinity is
\[
     |\mathrm{Sel1}_m|\le |\mathrm{Bd1}_m|Q^{md+C_1},
\]
where $C_1$ is the single constant from $B1$.

Do the same with the even and $W_2$ actions and the eventual
$B2$ statement with its own constant $C_2$. An array
$z_{ij}\in\mathrm{SelP}_m$ on the limit transports to that one
level by the even comparison. Apply $B2$ there. It gives a single
row $(c_i)_{0\le i\le g}$ over $A$, with a unit entry, and a
single $\Psi_N\in\mathrm{Sel2}_m$ whose $f_j$ coefficients
are $t^{C_2}\sum_i c_i z_{ij}$ on the transported array.
Use the $W_2$ comparison to transport $\Psi_N$ back. Equalities
of the specified coefficients hold on $P^0$, hence on the limit
by surjectivity. This proves the finite $B2$ conclusion at every
precision on the limit group, without changing $C_2$. The
ultrafilter levels may vary with precision; the constants do not.
% END L71C2

% BEGIN L71C3
\paragraph{L71C3. The two integral conclusions from the finite bounds.}
For $K1$ take $b$ integral $W_1$ Selmer cocycles whose $A$-linear
combinations are independent modulo the specified skew
coboundaries. Test whether a combination is a coboundary by the
linear equations in its $b$ coefficients and in all entries of
a primitive $Y$. Include the skew condition, the zero-coordinate
and block conditions on $Y$, and
$\sum_i c_i z_i(s)=R(s)YR(s)^{-1}-Y$ for all group elements.
These rows belong to one finite free dual $A$-module, since the
coefficient algebra is finite free and all these relations are
$A$-linear. The submodule they span is finitely generated.
For a solution modulo $t^m$, lift its variable entries to $A$;
every row then vanishes modulo $t^m$. Smith normal form for the
finite matrix of rows shows that the tuple agrees with an exact
kernel tuple modulo $t^{\max(m-b_0,0)}$, where $b_0$ is the
maximum valuation of a nonzero invariant factor (use zero for
an empty list). The coefficient part of every exact kernel
tuple is zero. Hence every coefficient of a relation modulo
boundaries at depth $m$ is divisible by $t^{\max(m-b_0,0)}$.
Boundaries lie in $\mathrm{Sel1}_m$: a skew primitive gives the
polarization identity and supplies its own local correction;
at tame inertia the deformed action is trivial. Reduction preserves
all Selmer conditions and all blocks. Distinct representatives
modulo this power therefore give distinct cosets of boundaries.
L71C2 yields $Q^{b\max(m-b_0,0)}\le Q^{md+C_1}$. Letting
$m$ grow proves $b\le d$, the integral $K1$.

For $K2$ fix an \emph{integral} even array and apply L71C2
to its reductions. At depth $m>0$ keep the row modulo $t^m$,
the whole continuous lifting cocycle in $W_2(A_m)$, and one
correcting matrix at each ordinary place. There are finitely
many such choices: the group at infinity is a quotient of $P$,
its sign kernel is topologically finitely generated, and the
argument with the finite affine group in L71C1 bounds all its
cocycles. The local matrices and rows are in finite sets.
Keep just the solutions with some unit coordinate and the
required $f_j$ equalities. These nonempty finite sets are
compatible under reduction to smaller positive depths. A compatible
sequence exists: at each depth keep the nodes having descendants
at arbitrarily large depths; among finitely many children at least
one has this property. Start at depth one and repeat.

Its rows and local matrices give $c_i,Y_v$ over $A$ by completeness.
Some $c_i$ is a unit, as seen at depth one. Pointwise inverse
limits give one cocycle $\Psi$ whose reductions are continuous
at each depth. Separatedness now gives its exact cocycle equations,
block conditions, transpose, tame and local equations, including
the same limiting $Y_v$ simultaneously for all local elements.
The $f_j$ coefficients equal $t^{C_2}\sum_i c_i z_{ij}$.
The row $(t^{C_2}c_i)_i$ is nonzero in the domain $A$, as
required by integral $K2$. This argument includes $g=0$.
% END L71C3

% BEGIN L71C4
\paragraph{L71C4. The integral odd classes descend through the actual limit kernel.}
Let $N$ be the closed normal subgroup of $\Gamma_\infty$
generated by its auxiliary elements $\tau_j$. Then
$N=\ker h$. One inclusion follows from $h(\tau_j)=1$.
For the converse, take a finite quotient $D$ of $\Gamma_\infty$
killing $N$, with map $\delta$. Write
$V=\ker(\delta\pi_\infty)$, an open normal subgroup of $P$.
At ultrafilter-many levels convergence of the kernels puts
$\ker\pi_N$ in $V$. Thus $\delta\pi_\infty$ factors through
a map $\delta_N:\Gamma_N\to D$ at those levels. Convergence of
the finitely many $\tau_{Nj}$ modulo $V$ forces
$\delta_N(\tau_{Nj})=1$ for all $j$. The kernel condition for
$T$ says that $\ker h_N$ is contained in the closed normal
subgroup generated by these elements and the relative
pro-$\ell$ residual. The latter is trivial by $T.\mathrm{ProL}$.
Consequently $\delta_N=\beta_Nh_N$ with continuous
$\beta_N:\Gamma_0\to D$. There are only finitely many such
homomorphisms, since $\Gamma_0$ is itself a continuous quotient
of the finitely generated $P$. On an ultrafilter-large subset
take $\beta_N=\beta$. Base projection convergence modulo
$\ker\beta$ gives
$\delta\pi_\infty=\beta h\pi_\infty$, hence $\delta=\beta h$.
All finite quotients of $\Gamma_\infty/N$ together separate
its points, proving the other inclusion.

Take any integral odd Selmer cocycle $z$ on the limiting sign
kernel. The elements $\varphi_j,\tau_j$ there commute,
$h(\tau_j)=1$, and $h(\varphi_j)=\gamma_j$ is the selected
base element. Comparing the cocycle on the two products implies
that $z(\tau_j)$ commutes with $r(\gamma_j)$. This matrix is
block diagonal, with separable characteristic polynomial split
in $E$ by the prescribed distinct roots. The eigenvalues of the
two blocks are disjoint; an intertwiner between them is zero.
Thus its odd centralizer over $E$ is zero, and so is its odd
centralizer over the subring $A$. We have $z(\tau_j)=0$.

The subgroup of $\ker h$ where $z$ vanishes is closed, since
$z$ is continuous at all precisions and $A$ is separated. It
is normal in the limiting sign kernel by the cocycle identity,
since $\ker h$ acts trivially on the coefficient module.
It is stable under the distinguished conjugation by the transpose
law, hence normal in the full group. The kernel calculation
puts all of $\ker h$ in it. Thus $z$ is the pullback of a
continuous base cocycle $z_0$. Local correction matrices can be
kept: the two frames have the same local source groups and
$h$ commutes with their maps. Similarly the inertia and
polarization conditions descend, so $z_0$ is in the original
odd Selmer module. The generator field of $k_x$ supplies $s,c_j$
and a skew odd primitive $Y$ such that
$t^s z_0-\sum_j c_j\kappa_j=dY$. Pulling back gives the
literal integral $K3$ on the limit. Together with L71C3 this
establishes $\mathrm{CohLim}$ for arbitrary block sizes and $g$.
% END L71C4

% BEGIN L71C5
\paragraph{L71C5. Passage to bounded Laurent cocycles and their local witnesses.}
In the source a bounded scalar Laurent function is one for which
$t^s$ times the function is an $A$-valued function continuous at
every $t$-adic precision, for some $s$. For a matrix or for a
matrix with entries in the paired coefficient algebra this
condition holds coordinate by coordinate, including its constant,
$e_j,f_j,\eta$ coordinates. For any finite family of such
cocycles take a common exponent, large enough also to clear the
denominators of the entries of a fixed correcting matrix at
each local place for each cocycle. The local set and all matrix
index sets are finite. Scaling each witness and each cocycle by
this same power makes them integral. The action and change of
flag are integral, their equations linear in the cocycle and
witness, and all block and coordinate restrictions are zero
equations. Injectivity into $E$ therefore proves all these
identities integrally, including the torsion diagonal restrictions.
Continuity of the integral entries follows from the boundedness
witnesses. Conversely integral cocycles extend with bounded
coordinates, by taking exponent zero.

For a family as in the Laurent rank bound, scale it by a common
$t^s$. An integral relation among its scaled members modulo
integral skew boundaries gives a Laurent relation among the
original family by applying the coefficient inclusion. The
hypothesized Laurent independence makes every coefficient zero
since $t^s\ne0$. Apply L71C3 to bound the size by $d$.
This gives $\mathrm{SelmerRankBoundLaurent}$.

For the rectangular array of bounded even cocycles, scale again
by one $t^s$ before using integral $K2$. Its lifting cocycle
extends with bounded Laurent coefficients. The resulting
nonzero Laurent row has coordinates $t^s c_i$ if $(c_i)$ is the
nonzero integral row for the scaled array. Thus its ordered
$f_j$ components are exactly the asserted combination of the
original unscaled rows. This gives $\mathrm{EvenCocyclesLiftLaurent}$.
For an odd bounded $z$, integral $K3$ gives
$t^{b+s}z-\sum_j c_j\kappa_j=dY$ after coefficient inclusion,
where $t^s z$ was the integral input. Divide the coefficients
and primitive by $t^{b+s}$ in $E$, preserving the skew odd
condition, to obtain $\mathrm{SelmerSpannedByKappaLaurent}$.
These three statements hold on the same point, package, projection
and limit used in R70P5. They require only the stated original
$B1$, $B2$, $T.\mathrm{ProL}$, auxiliary eigenvalues
and the limit supplied by LIM18.
% END L71C5
